% ==============================================================================
% reports/sections/discussion/07_01_zero_shot_challenge.tex
% ==============================================================================

\section{Tại Sao Bài Toán Phân Tách Theo Thời Gian Lại Khó: Thách Thức Subtype Distribution Shift}

Sự sụp đổ của cả 3 mô hình học máy dạng cây trên tập Test theo thời gian ($\FOne \le \RFTestFOne$) là phát hiện khoa học mang tính thực tiễn cao nhất của đồ án này. Trong khi các nghiên cứu trước đây thường che giấu hoặc tránh né kịch bản này, chúng tôi khẳng định đây là một hiện tượng tự nhiên và tất yếu bắt nguồn từ \textbf{Sự dịch chuyển phân phối của các biến thể tấn công (Subtype Distribution Shift)}:

\subsection{Sự Khác Biệt Căn Bản Về Cấu Trúc Giao Thức Mạng}
FTP và SSH đại diện cho hai triết lý thiết kế giao thức mạng hoàn toàn trái ngược nhau:
\begin{enumerate}
    \item \textbf{FTP (Truyền thông văn bản rõ)}: Quá trình bắt tay diễn ra tối giản. Các gói tin yêu cầu \texttt{USER} và \texttt{PASS} có độ dài rất ngắn và kích thước cố định. Do đó, các đặc trưng như \texttt{Fwd Packet Length Std}, \texttt{Fwd Packet Length Mean} và \texttt{Average Packet Size} tập trung ở các dải giá trị rất hẹp và thấp.
    \item \textbf{SSH (Truyền thông mã hóa kênh)}: Trước khi bước vào giai đoạn thử mật khẩu, SSH bắt buộc phải thực hiện pha trao đổi khóa mật mã phức tạp (Diffie-Hellman Key Exchange) và thương lượng thuật toán mã hóa đối xứng. Pha trao đổi này làm phát sinh các gói tin có kích thước lớn hơn gấp nhiều lần so với FTP, dẫn đến phân phối của \texttt{Packet Length Mean} và \texttt{Max Packet Length} dịch chuyển hoàn toàn sang một miền giá trị mới.
\end{enumerate}

\subsection{Bản Chất Của Thuật Toán Học Máy Giám Sát Cổ Điển}
Các mô hình học máy dạng cây (Decision Tree, Random Forest, XGBoost) hoạt động dựa trên nguyên lý \textbf{học phân vùng không gian đặc trưng có giám sát (Supervised Partitioning)}. Khi toàn bộ dữ liệu huấn luyện $\Dtrain$ chỉ bao gồm các mẫu tấn công FTP, không gian quyết định (Decision Boundary) của mô hình bị khóa chặt vào vùng đặc trưng đại diện cho FTP:
\begin{equation}
    \mathcal{R}_{\text{Attack}} = \left\{ \mathbf{x} \mid \text{Port} \le 21.5 \land \text{Fwd Packet Length Std} \le \theta_{\text{FTP}} \dots \right\}
\end{equation}

Khi một cuộc tấn công SSH xuất hiện, tọa độ của nó trong không gian đặc trưng rơi hoàn toàn vào vùng không gian chưa từng được gán nhãn tấn công trong quá trình huấn luyện:
\begin{equation}
    \mathbf{x}_{\text{SSH}} \notin \mathcal{R}_{\text{Attack}} \implies P(y=1|\mathbf{x}_{\text{SSH}}) \approx 0
\end{equation}

Mô hình học máy không sở hữu khả năng suy luận logic bậc cao như con người để hiểu rằng "cả hai đều là hành vi dò quét mật khẩu lặp đi lặp lại". Đối với mô hình, SSH đơn thuần là một phân phối thống kê hoàn toàn xa lạ, và do đó nó bị phân loại nhầm thành lưu lượng bình thường (\Normal) với độ tin cậy tuyệt đối.
